%##########################################################################
% CHAPTER 6: DISCUSSION
%##########################################################################
\chapter{Discussion}

\section{Chapter Overview}

This chapter steps back from the numbers and themes presented in Chapter~5 and asks what they mean. It interprets the quantitative and qualitative findings in relation to the research questions and the existing literature, identifies the original contributions of the work, and --- with equal candour --- acknowledges where the system falls short and why.

\section{Interpretation of Quantitative Findings}

The central prediction of this research --- that PPO-driven narrative adaptation would produce measurably higher Flow experience --- rests on a theoretical chain with several links. The player modelling module must accurately infer the player's psychological state from noisy behavioural telemetry. The PPO agent must select a narratively appropriate action in response to that inferred state. The LLM must generate content that faithfully executes the selected action while remaining grounded in the game's lore. And the player must actually experience the resulting narrative shift as meaningful --- not as noise, not as interruption, but as a story that responds to them. If any link in this chain fails, the effect on the dependent variable (GEQ Flow subscale) will be attenuated or absent.

\textit{[INSERT: Interpret the actual results when available. The following structures illustrate possible interpretations:]}

\textbf{If H1 is supported:} The significant difference in GEQ Flow scores ($d$ = \textit{[INSERT]}) supports the core theoretical claim that real-time narrative personalization can improve Flow experience. The effect size is \textit{[magnitude]} by Cohen's (\citeyear{cohen1988statistical}) convention, suggesting that narrative adaptation is a meaningful lever for engagement even within a relatively short (30-minute) session. The convergent evidence from the behavioural metrics --- specifically the dialogue read ratio in the Experimental condition --- suggests that players were not merely reporting higher Flow retrospectively but were also demonstrably more engaged with the narrative content during play.

\textbf{If H1 is not supported:} The absence of a significant difference challenges the assumption that within-session narrative adaptation produces observable Flow differences within a 30-minute session. Several explanations merit consideration. First, the PPO agent was trained on a simulated environment whose transition dynamics only approximately match real player behaviour; policy transfer to the real game environment may be limited. Second, the 30-minute session may be insufficient for the player modelling module to accumulate enough telemetry. Third, the effect of narrative tone on Flow may be smaller than anticipated relative to mechanical game design factors.

\textit{[INSERT: Discuss H2 (Immersion) and H3 (Personalization) results in relation to each other and to H1.]}

\section{Themes from Interviews}

\textit{[INSERT: Discuss qualitative themes in relation to the quantitative results and existing literature.]}

The interview data are expected to illuminate the \textit{mechanisms} by which narrative adaptation influenced (or failed to influence) experience --- something the quantitative measures, by design, cannot reveal. A structured scale tells you \textit{whether} a difference exists; a semi-structured interview tells you \textit{why} it exists, and --- when the difference is absent --- \textit{what went wrong}.

A theme of \textbf{narrative responsiveness} --- players describing the story as attuned to their behaviour, as if the game were ``paying attention'' --- would provide direct qualitative support for the personalization hypothesis (H3) and would suggest that the adaptation pipeline's effect is perceptible at the level of subjective experience, not merely at the level of aggregated scale scores. Conversely, a theme of \textbf{artificiality awareness} --- players noticing repetitive patterns, inconsistent character voices, anachronistic language, or content that does not fit the current situation --- would identify specific failure modes of the LLM generation pipeline. If multiple participants independently describe the same kind of failure (e.g., ``the dialogue felt generic'' or ``the tone didn't match what was happening''), that convergence would provide actionable feedback for improving the prompt engineering and content validation stages of the system.

The Flow-related interview questions are expected to reveal whether narrative injections were perceived as supportive of engagement or as disruptive interruptions. This distinction matters greatly. A recurring finding in studies of game storytelling is that narrative interruptions during high-engagement moments --- a cutscene triggered in the middle of an intense combat encounter, for example --- are perceived negatively even when the content itself is high quality. The \texttt{ContentInjector}'s 3-second minimum spacing and its COOLING\_DOWN state were designed to mitigate this risk, but only the interview data can reveal whether the mitigation was sufficient in practice.

\section{Novelty and Contributions}

The original contributions of this work are six in number, and they range from the practical (a reusable mapping table) to the conceptual (reframing what reinforcement learning in games is for).

\textbf{Contribution 1: Telemetry-derived Hexad profiling.} Every prior system that has used Hexad player types for game adaptation has relied on a pre-play questionnaire \citep{tondello2016hexad} --- a static snapshot that cannot update during a session. This system computes six-dimensional continuous profiles directly from behavioural telemetry, enabling within-session updating without any self-report instrument. The telemetry-to-Hexad mapping (Table~\ref{tab:hexad}) is published in full and is reusable by other researchers and developers working with different games and different telemetry schemas.

\textbf{Contribution 2: RAG-grounded dynamic narrative.} PANGeA \citep{todd2024pangea} grounds LLM narrative in prompt-delivered context, but that context is assembled at design time and remains fixed across players and sessions. In the present system, the lore passages delivered to the LLM are selected at runtime based on the player's current location, activity, and the narrative action chosen by the PPO agent. The grounding is therefore dynamic rather than static: the same player in the same location at different Flow states will receive different lore context, producing different narrative output.

\textbf{Contribution 3: Episodic session memory in the LLM prompt.} The importance-weighted episodic memory, inspired by EM-LLM \citep{fountas2024episodic}, gives the LLM access to a compressed narrative history so that generated content can reference prior events, maintain character continuity, and avoid contradicting what the player has already experienced. To the author's knowledge, this is the first application of episodic memory compression to real-time game narrative generation.

\textbf{Contribution 4: PPO for narrative action selection.} Reinforcement learning has been applied extensively to mechanical game adaptation --- adjusting enemy health, spawn rates, and resource drops --- but not, to the author's knowledge, to selecting \textit{narrative actions}. Framing narrative tone selection as an RL problem is a conceptual step beyond both rule-based narrative adaptation and mechanical DDA. It implies that the ``right'' narrative response to a given player state is not something a designer can fully specify in advance; it must be learned from the interaction dynamics.

\textbf{Contribution 5: Flow-based MDP reward.} The reward function directly operationalises Csikszentmihalyi's (\citeyear{csikszentmihalyi1990flow}) challenge--skill balance theory, giving the adaptation objective a principled theoretical grounding rather than an ad hoc engagement metric. The graduated penalty structure (APATHY $>$ ANXIETY $>$ BOREDOM) reflects the severity ordering of disengagement states, and the shaping rewards (streak bonus, transition bonus, repetition penalty) encode design-level knowledge about what good adaptation looks like.

\textbf{Contribution 6: Complete closed-loop integration.} The five modules --- telemetry capture, feature extraction and player typing, RL action selection, RAG-grounded LLM generation, and in-game content injection --- run end-to-end within a single deployable Unity plugin. Individual pieces of this pipeline have appeared in prior work; what has not appeared is the closed loop. The player's behaviour influences the narrative, and the narrative influences the player's behaviour, creating a feedback cycle that is the defining characteristic of a genuinely adaptive system.

\section{Limitations}

No system is stronger than its weakest assumption, and this one has several that warrant explicit acknowledgement.

\textbf{Simulation fidelity.} The PPO agent is trained on a simulated environment with hand-specified transition probabilities --- a 0.75 probability that a ``helpful'' action will move the player to FLOW, and a 0.25 probability that it will not. Real player behaviour is substantially more complex, more stochastic, and more history-dependent. The simulation's Markovian assumption --- that the next state depends only on the current state and action, not on the sequence of states that preceded it --- is almost certainly violated in practice. A player who has been in ANXIETY for five minutes is not in the same psychological position as one who just entered ANXIETY from FLOW, even if the telemetry-derived state vector is identical. The trained policy may therefore underperform in the real game environment relative to its performance in simulation.

\textbf{Rule-based Hexad profiling.} The mapping from telemetry signals to Hexad dimensions is based on theoretical alignment with Tondello et al.'s (\citeyear{tondello2016hexad}) motivation descriptions, not on empirical validation against questionnaire-derived ground truth. This means that the system's internal Hexad profiles may diverge from what a validated questionnaire would produce for the same player. Systematic misclassification is possible: a player who moves slowly due to hardware input lag or unfamiliarity with the controls, rather than deliberate exploration, will register a high \texttt{explore\_rate} and be profiled as an Explorer. The system has no way to distinguish intentional from incidental behaviour.

\textbf{LLM generation quality.} Phi-3.5 Mini is a 3.8-billion-parameter model --- capable for its size, but fundamentally constrained relative to larger frontier models. The quality of generated narrative --- its fluency, creativity, emotional range, and consistency of character voice --- is bounded by the model's capacity. Larger models would likely produce more compelling output, but they would also introduce latency, memory, and deployment constraints that conflict with the real-time requirement.

\textbf{Single-session study.} The evaluation captures a single 30-minute session per participant. Any effects of narrative personalization that require multiple sessions to develop --- deepening attachment to characters, evolving expectations, long-term engagement --- are invisible to this design. A longitudinal study would be needed to assess cumulative impact.

\textbf{Testbed ecological validity.} The testbed game is a simplified prototype built on a generic RPG framework. It lacks the production polish, narrative complexity, and mechanical depth of a commercially developed title. Findings obtained in this controlled environment may not generalise to games with more complex narrative architectures, larger worlds, or longer play sessions.

\textbf{Sample size.} With $n = 25$ per condition, the study is powered to detect medium-to-large effects ($d \geq 0.5$) but would miss small effects entirely. The convenience sample --- drawn from university students and gaming society members --- also limits demographic diversity and may not represent the broader population of RPG players.

\section{Threats to Validity}

\textbf{Internal validity.} The primary threat is demand characteristics. Participants who infer that they are in the experimental condition --- perhaps because they notice that the narrative seems to be responding to their behaviour --- may respond more favourably on the post-session questionnaires, inflating the apparent effect. The blind design (participants are not told which condition they are in) mitigates this threat, but does not eliminate it entirely. A more robust mitigation would be a double-blind design in which the researcher administering the session is also unaware of condition assignment, but the single-researcher study setup makes this impractical.

\textbf{Construct validity.} The GEQ Flow subscale measures retrospective self-report of Flow --- what the participant \textit{remembers} feeling during the session, not what they \textit{actually} felt moment to moment. Peak-end effects and recency bias may distort the score. The absence of a physiological measure (e.g., galvanic skin response, heart rate variability) limits confidence that the GEQ score accurately reflects the within-session psychological state. The custom NPS-3 scale, while face-valid, has not been independently validated beyond the present study.

\textbf{External validity.} The testbed game is a simplified prototype, and the participants are a convenience sample of university students. Both factors limit generalisability: the system's effectiveness in a more complex commercial game with a more diverse player base remains an open question.

\textbf{Statistical conclusion validity.} With $N = 50$ and Bonferroni correction across three hypotheses ($\alpha_{\text{adj}} = 0.0167$), the study is not robustly powered for small effects. A true but small effect ($d < 0.4$) would likely fail to reach significance, leading to a Type II error. This limitation is acknowledged in the design and motivates the inclusion of qualitative data as a complementary source of evidence.
